% E-from-v1.tex: what Version 2 kept from Version 1, and what it corrected.
\section{From Version 1 to Version 2}
\label{app:from-v1}

Version 1 remains the historical report and is not edited. This appendix records what Version 2
kept, what it corrected, and what it retired, so that a reader of both can reconcile them.

\subsection*{Kept}

\begin{itemize}
  \item The token-centred entity geometry (Figure~\ref{fig:nucleus}), reproduced from Version 1
    without change and used as the nucleus of every map in this paper.
  \item The lifecycle state machine (Figure~\ref{fig:lifecycle}), reproduced without change.
  \item The disclosure-level argument: that the level is a persisted property of the verification
    record, enforced by a check constraint, and that this is what prevents the verification log from
    functioning as a surveillance database.
  \item The argument that post-quantum signing belongs on the first day, through the secrecy
    horizon of identity data and Mosca's inequality, and the substrate argument that every higher
    property of an identity system is derivative of its cryptographic primitives. Both are summarised
    here and stand in Version 1 in full.
  \item The comparison of design properties against deployed systems, which Version 2 does not
    repeat; the repository's front page carries the current table and marks Polaris as neither
    deployed nor issuing at national scale.
  \item The threat-model frame (STRIDE), which the repository's design record now carries with every
    threat mapped to a control or an explicit acceptance.
\end{itemize}

\subsection*{Corrected}

\begin{itemize}
  \item Twelve tables became thirty-six. The three subsidiary entities of Version 1 (device
    bindings, ledger anchors, the revocation list) remain; the rest of the growth is described ring by
    ring in Section~\ref{sec:walls}.
  \item Version 1 stated that the audit tables were append-only by convention and tooling and that
    the state machine would be enforced by a trigger in a production deployment. Both are now
    enforced by triggers, and a check fails the build if either stops being.
  \item Version 1 described the ledger anchor as an optional decentralised-identifier path. The
    off-chain batch commitment exists and is monitored as a transparency log; publication to an
    external chain remains an operator action with declared drivers and no live integration.
  \item Version 1 modelled seven verification contexts with per-context biometric requirements and a
    physical token with on-device biometric binding. The contexts remain; the physical token remains
    modelled and not manufactured, and this paper marks it as a proposal.
  \item The two speculative binding horizons of Version 1's final appendix (genomic and
    quantum-observer binding) are not carried forward as design. The repository's readiness pass
    quarantined the speculative schema behind a check that it cannot return, and this paper treats
    the holder-side key, not a new biometric, as the extension that the future environment needs.
\end{itemize}

\subsection*{Added}

Everything from Section~\ref{sec:signatures} onward describes machinery that did not exist when
Version 1 was written: real signing and custody, the two witnesses, the detached verifier and the
conformance contract, status and offline authorization, the membership proof and its second witness,
the wallet protocol and the broker, federation, transparency, the exchange fabric, the operational
drills, Atlas and Athena, the history of self-correction, the societal boundaries, and the
future-facing sections on agents, personhood, delegation and outside observers.
